\addcontentsline{toc}{subsubsection}{Solution 42. Bounded Errors, Exact Answer}

\begin{dlsummary}[title={\bfseries Solution 42. Bounded Errors, Exact Answer},breakable]

\begin{proof}
The only solution is
\[
\boxed{f(x)=x\qquad(x\in\mathbb R).}
\]

\paragraph{1. Force unboundedness.}

Repeatedly applying the additive inequality gives
\[
f(2n)\ge nf(2)-(n-1)
=1+n\bigl(f(2)-1\bigr).
\]
Since \(f(2)>1\), these values tend to infinity.

\paragraph{2. Eliminate the multiplicative error.}

Write
\[
E(x,y)=f(xy)-f(x)f(y),\qquad |E(x,y)|\le1.
\]
Evaluating the same triple product in two ways yields
\[
E(x,y)f(t)
=E(x,yt)-E(xy,t)+f(x)E(y,t).
\]
Consequently,
\[
|E(x,y)|\,|f(t)|\le2+|f(x)|.
\]
For fixed \(x,y\), the right-hand side is constant, while \(|f(t)|\) can be arbitrarily large. Therefore \(E(x,y)=0\), proving
\[
f(xy)=f(x)f(y). \tag{1}
\]

\paragraph{3. Eliminate the additive error.}

Apply the additive inequality to \(tx,ty\). Using (1),
\[
|f(t)|\,|f(x+y)-f(x)-f(y)|\le1.
\]
Again, unboundedness forces
\[
f(x+y)=f(x)+f(y). \tag{2}
\]

\paragraph{4. Recover order, then the function.}

Multiplicativity gives \(f(1)^2=f(1)\); since \(f\) is not identically zero, \(f(1)=1\). Additivity now gives
\[
f(q)=q\qquad(q\in\mathbb Q).
\]

For \(u>0\),
\[
f(u)=f(\sqrt u)^2>0,
\]
where strict positivity follows from
\[
f(\sqrt u)\,f(1/\sqrt u)=1.
\]
Thus, if \(b>a\), then
\[
f(b)-f(a)=f(b-a)>0.
\]
So \(f\) is strictly increasing. Comparing any real \(x\) with rational numbers immediately below and above it forces \(f(x)=x\).

The identity function plainly satisfies both inequalities and \(f(2)>1\).
\end{proof}
\end{dlsummary}

\begin{takeaways}
\begin{itemize}
\item \textbf{Associativity can remove errors:} Comparing two expansions of a triple product makes a fixed discrepancy multiply an unbounded quantity.
\item \textbf{Order can be a conclusion:} Squares establish positivity, and additivity turns positivity into monotonicity.
\item \textbf{A threshold can matter:} The condition \(f(2)>1\) excludes \(f\equiv1\), which satisfies both error bounds.
\item \textbf{Avoid assuming regularity:} The argument establishes every property it needs directly from the hypotheses.
\end{itemize}
\end{takeaways}
